% =====================================================================
\chapter{Teorema fundamental del cálculo, áreas y valor medio}
\label{tema:2}
% =====================================================================

El \cref{tema:1} definió la integral como un límite de sumas, una definición rigurosa pero poco práctica: calcular $\int_0^1x^{3}\dif x$ por sumas ya exige conocer $\sum k^{3}$. El teorema fundamental del cálculo cambia por completo la situación, porque reduce el cálculo de una integral al de una primitiva y, a la vez, garantiza que toda función continua tiene primitiva. Este tema lo demuestra en sus dos partes y lo aplica a áreas, valores medios y funciones definidas por integrales, como el logaritmo. Se siguen el capítulo~14 de Spivak~\cite{spivak}, el capítulo~5 de Apostol~\cite{apostol}, la §5.3 de~\cite{openstax1}, las §5.6 y~5.7 de Strang~\cite{strang} y las partes~A a~C de la unidad~3 del curso 18.01SC~\cite{mit1801}.

\section{El teorema del valor medio integral}

\begin{teorema}[Valor medio integral]\label{teo:tvm-integral}
Sea $f$ continua en $[a,b]$. Existe $c\in[a,b]$ tal que
\[
\int_a^b f(x)\dif x=f(c)\,(b-a).
\]
\end{teorema}

\begin{proof}
Por el teorema de Weierstrass (módulo~0.5), $f$ alcanza un mínimo $m=f(x_1)$ y un máximo $M=f(x_2)$ en $[a,b]$. Por el \cref{teo:propiedades}(iv), $m(b-a)\le\int_a^b f\le M(b-a)$, es decir, el número $\mu=\frac{1}{b-a}\int_a^b f$ cumple $f(x_1)\le\mu\le f(x_2)$. Por el teorema de los valores intermedios aplicado a $f$ en el intervalo de extremos $x_1$ y $x_2$, existe $c$ entre ellos con $f(c)=\mu$.
\end{proof}

\begin{definicion}[Valor medio de una función]\label{def:valor-medio}
Sea $f$ integrable en $[a,b]$. Su \emph{valor medio} en $[a,b]$ es
$\displaystyle\bar f=\frac{1}{b-a}\int_a^b f(x)\dif x$.
\end{definicion}

Geométricamente, si $f\ge0$, el rectángulo de base $[a,b]$ y altura $\bar f$ tiene la misma área que la región bajo la gráfica. El \cref{teo:tvm-integral} dice que una función continua toma su valor medio en algún punto. El nombre procede de la analogía con la media de $n$ números: $\bar f$ es el límite de las medias $\frac1n\sum_{k=1}^{n}f(x_k)$ en puntos equiespaciados, que son sumas de Riemann divididas por $b-a$.

\section{El teorema fundamental del cálculo}

\begin{teorema}[Teorema fundamental del cálculo, primera parte]\label{teo:tfc1}
Sea $f$ integrable en $[a,b]$ y sea $F(x)=\int_a^x f(t)\dif t$ para $x\in[a,b]$.
\begin{enumerate}
  \item $F$ es continua en $[a,b]$; de hecho, si $\abs{f}\le K$, entonces $\abs{F(y)-F(x)}\le K\abs{y-x}$.
  \item Si $f$ es continua en un punto $c\in(a,b)$, entonces $F$ es derivable en $c$ y $F'(c)=f(c)$.
\end{enumerate}
\end{teorema}

\begin{proof}
(i) Sea $x<y$. Por la aditividad, $F(y)-F(x)=\int_x^y f$, y por la acotación del \cref{teo:propiedades}, $\abs{F(y)-F(x)}\le K(y-x)$.

(ii) Sea $h\neq0$ con $c+h\in[a,b]$. Como $\int_c^{c+h}f(c)\dif t=f(c)\,h$ (también si $h<0$, por el convenio de la \cref{def:integral}),
\[
\frac{F(c+h)-F(c)}{h}-f(c)=\frac1h\int_c^{c+h}\bigl(f(t)-f(c)\bigr)\dif t .
\]
Sea $\varepsilon>0$. Por la continuidad de $f$ en $c$ existe $\delta>0$ tal que $\abs{f(t)-f(c)}<\varepsilon$ si $\abs{t-c}<\delta$. Si $0<\abs{h}<\delta$, todos los $t$ entre $c$ y $c+h$ cumplen $\abs{t-c}<\delta$, y la acotación da $\bigl|\int_c^{c+h}(f(t)-f(c))\dif t\bigr|\le\varepsilon\abs{h}$. Dividiendo por $\abs{h}$, el cociente incremental dista de $f(c)$ menos de $\varepsilon$ (o a lo sumo $\varepsilon$). Luego $F'(c)=f(c)$.
\end{proof}

\begin{corolario}\label{cor:existencia-primitiva}
Toda función continua en un intervalo $I$ tiene primitiva en $I$: para cualquier $a\in I$, $F(x)=\int_a^x f(t)\dif t$ lo es.
\end{corolario}

\begin{corolario}[Derivación de integrales con límites variables]\label{cor:leibniz}
Sea $f$ continua en un intervalo $I$ y sean $g,h$ derivables con valores en $I$. Entonces
\[
\frac{d}{dx}\int_{h(x)}^{g(x)}f(t)\dif t=f\bigl(g(x)\bigr)\,g'(x)-f\bigl(h(x)\bigr)\,h'(x).
\]
\end{corolario}

\begin{proof}
Sea $a\in I$ y $F(u)=\int_a^u f$. Por la aditividad, $\int_{h(x)}^{g(x)}f=F(g(x))-F(h(x))$. Por el \cref{teo:tfc1}, $F'=f$, y la regla de la cadena da el resultado.
\end{proof}

\begin{teorema}[Teorema fundamental del cálculo, segunda parte, o regla de Barrow]\label{teo:barrow}
Sea $f$ integrable en $[a,b]$ y sea $F$ continua en $[a,b]$, derivable en $(a,b)$, con $F'(x)=f(x)$ para todo $x\in(a,b)$. Entonces
\[
\int_a^b f(x)\dif x=F(b)-F(a)=\eval{F(x)}{a}{b}.
\]
\end{teorema}

\begin{proof}
Sea $P=\{x_0,\dots,x_n\}$ una partición cualquiera. La suma telescópica da $F(b)-F(a)=\sum_{i=1}^{n}\bigl(F(x_i)-F(x_{i-1})\bigr)$. En cada $[x_{i-1},x_i]$, el teorema del valor medio de Lagrange proporciona $t_i\in(x_{i-1},x_i)$ con $F(x_i)-F(x_{i-1})=F'(t_i)\Delta x_i=f(t_i)\Delta x_i$. Como $m_i\le f(t_i)\le M_i$,
\[
\Lsum(f,P)\le F(b)-F(a)=\sum_{i=1}^{n}f(t_i)\Delta x_i\le\Usum(f,P).
\]
El número $F(b)-F(a)$ está, pues, entre todas las sumas inferiores y todas las superiores; como $f$ es integrable, el único número con esa propiedad es $\int_a^b f$.
\end{proof}

La demostración no exige que $f$ sea continua, solo integrable, y no usa la primera parte. Cuando $f$ es continua también puede deducirse de ella: $G(x)=\int_a^x f$ y $F$ son primitivas de $f$ en $(a,b)$, continuas en $[a,b]$, luego difieren en una constante (\cref{prop:primitivas-constante}), y $F(b)-F(a)=G(b)-G(a)=\int_a^b f$.

\begin{ejemplo}\label{ej:barrow}
(a) $\int_0^{\pi}\sin x\dif x=\eval{-\cos x}{0}{\pi}=1-(-1)=2$. (b) $\int_0^1x^{n}\dif x=\eval{\frac{x^{n+1}}{n+1}}{0}{1}=\frac{1}{n+1}$, lo que, con $n=2$, confirma el \cref{ej:x2} sin sumas. (c) $\int_0^1\frac{\dif x}{1+x^{2}}=\arctan1-\arctan0=\frac\pi4$.
\end{ejemplo}

\section{El logaritmo como integral}

El \cref{cor:existencia-primitiva} permite \emph{definir} funciones mediante integrales. El ejemplo más importante es el logaritmo, que puede construirse sin conocer previamente la exponencial~\cite[cap.~18]{spivak}, \cite[cap.~6]{apostol}.

\begin{proposicion}\label{prop:logaritmo}
Sea $L(x)=\int_1^x\frac{\dif t}{t}$ para $x>0$. Entonces $L(1)=0$, $L'(x)=\frac1x$, $L$ es estrictamente creciente y $L(ab)=L(a)+L(b)$ para todos $a,b>0$.
\end{proposicion}

\begin{proof}
$L(1)=\int_1^1=0$. Como $t\mapsto1/t$ es continua en $(0,\infty)$, el \cref{teo:tfc1} da $L'(x)=1/x>0$, y $L$ es estrictamente creciente. Fijado $a>0$, sea $g(x)=L(ax)-L(x)$ en $(0,\infty)$. Por la regla de la cadena, $g'(x)=a\cdot\frac{1}{ax}-\frac1x=0$, luego $g$ es constante (\cref{prop:primitivas-constante}) e igual a $g(1)=L(a)-L(1)=L(a)$. Con $x=b$: $L(ab)-L(b)=L(a)$.
\end{proof}

$L$ es, por tanto, el logaritmo natural, $L=\ln$, y su inversa es la exponencial. De la misma forma se definen funciones que no tienen primitiva expresable con funciones elementales, como la función de error $\operatorname{erf}(x)=\frac{2}{\sqrt\pi}\int_0^x e^{-t^{2}}\dif t$ o la función de distribución de la normal estándar, $\Phi(x)=\frac{1}{\sqrt{2\pi}}\int_{-\infty}^{x}e^{-t^{2}/2}\dif t$, que es una integral impropia (\cref{tema:4}). Por el \cref{teo:tfc1}, $\Phi'(x)=\frac{1}{\sqrt{2\pi}}e^{-x^{2}/2}$ aunque $\Phi$ no tenga fórmula elemental.

\section{Áreas entre curvas}

\begin{definicion}[Área entre dos curvas]\label{def:area}
Sean $f,g$ continuas en $[a,b]$. El área de la región comprendida entre sus gráficas y las rectas $x=a$, $x=b$ es $\int_a^b\abs{f(x)-g(x)}\dif x$. Si $f\ge g$ en $[a,b]$, es $\int_a^b(f-g)$.
\end{definicion}

Para calcularla se hallan los puntos donde $f-g$ cambia de signo, se divide $[a,b]$ en subintervalos donde el signo es constante y se suman las integrales de $f-g$ o de $g-f$ según corresponda.

\begin{ejemplo}\label{ej:area-parabola}
La recta $y=x+2$ y la parábola $y=x^{2}$ se cortan donde $x^{2}-x-2=0$, en $x=-1$ y $x=2$. En $[-1,2]$ la recta está por encima (en $x=0$: $2>0$). El área encerrada es
\[
\int_{-1}^{2}(x+2-x^{2})\dif x=\Bigl[\frac{x^{2}}{2}+2x-\frac{x^{3}}{3}\Bigr]_{-1}^{2}=\frac{10}{3}-\Bigl(-\frac76\Bigr)=\frac92 .
\]
\end{ejemplo}

\begin{ejemplo}\label{ej:area-seno}
El área entre $y=\sin x$ y el eje $x$ en $[0,2\pi]$ es $\int_0^{2\pi}\abs{\sin x}\dif x=\int_0^{\pi}\sin x\dif x-\int_{\pi}^{2\pi}\sin x\dif x=2-(-2)=4$, mientras que $\int_0^{2\pi}\sin x\dif x=0$.
\end{ejemplo}

\section{Valor medio y simetrías}

\begin{ejemplo}\label{ej:media-seno}
El valor medio de $\sin$ en $[0,\pi]$ es $\frac1\pi\int_0^{\pi}\sin x\dif x=\frac2\pi\approx0{,}6366$. El \cref{teo:tvm-integral} garantiza un $c$ con $\sin c=\frac2\pi$; hay dos, $c=\arcsin\frac2\pi\approx0{,}6901$ y $\pi-c$.
\end{ejemplo}

\begin{proposicion}[Integrales de funciones pares e impares]\label{prop:par-impar}
Sea $f$ continua en $[-a,a]$. Si $f$ es par, $\int_{-a}^{a}f=2\int_0^{a}f$; si $f$ es impar, $\int_{-a}^{a}f=0$.
\end{proposicion}

\begin{proof}
Sea $F(x)=\int_0^xf$ y sea $G(x)=\int_{-x}^{x}f=F(x)-F(-x)$ para $x\in[0,a]$. Por el \cref{cor:leibniz}, $G'(x)=f(x)+f(-x)$. Si $f$ es impar, $G'=0$ y $G(x)=G(0)=0$. Si $f$ es par, $G'(x)=2f(x)=\frac{d}{dx}\bigl(2F(x)\bigr)$, y como $G(0)=2F(0)=0$, $G=2F$.
\end{proof}

\begin{observacion}[Medias ponderadas]
Si $w\ge0$ es continua con $\int_a^b w=1$, el número $\int_a^b f(x)\,w(x)\dif x$ es una media de $f$ ponderada por $w$. Cuando $w$ es la densidad de una variable aleatoria $X$, es la esperanza $\E[f(X)]$ del módulo~5; el valor medio de la \cref{def:valor-medio} corresponde a la densidad uniforme $w=\frac{1}{b-a}$. El \cref{eje:2-tvm-ponderado} generaliza el \cref{teo:tvm-integral} a este caso.
\end{observacion}

\begin{observacion}[Errores frecuentes]\label{obs:2-errores}
(a) Olvidar el factor de la regla de la cadena: $\frac{d}{dx}\int_0^{x^{2}}f=f(x^{2})\cdot2x$, no $f(x^{2})$. (b) Aplicar la regla de Barrow cuando $f$ no es integrable o $F$ no es primitiva en todo el intervalo: «$\int_{-1}^{1}\frac{\dif x}{x^{2}}=\eval{-\frac1x}{-1}{1}=-2$» es falso, porque el integrando, positivo, no está acotado cerca de~$0$ (\cref{eje:2-error-barrow}). (c) Calcular un área como $\int_a^b(f-g)$ sin comprobar que las curvas no se cruzan. (d) Olvidar dividir por $b-a$ al calcular un valor medio.
\end{observacion}

\begin{observacion}[Conexión con el curso]
En el módulo~5, la función de distribución de una variable continua es $F_X(x)=\int_{-\infty}^{x}f_X(t)\dif t$; el \cref{teo:tfc1} da $F_X'=f_X$ en los puntos de continuidad de la densidad, y la regla de Barrow, $P(a\le X\le b)=F_X(b)-F_X(a)$. En el módulo~9, la activación $\mathrm{GELU}(t)=t\,\Phi(t)$ se deriva con el \cref{teo:tfc1}: $\mathrm{GELU}'(t)=\Phi(t)+t\,\Phi'(t)$, aunque $\Phi$ no tenga expresión elemental. En el módulo~7, la regla de Barrow aplicada a $\varphi(t)=f(x+td)$ en $[0,1]$ da la fórmula $f(x+d)-f(x)=\int_0^1\varphi'(t)\dif t$, base del lema de descenso.
\end{observacion}

\begin{lecturas}
\item Spivak~\cite{spivak}, cap.~14 «The Fundamental Theorem of Calculus».
\item Apostol~\cite{apostol}, cap.~5 «The Relation between Integration and Differentiation».
\item OpenStax~\cite{openstax1}, §5.3 «The Fundamental Theorem of Calculus» y §5.4 «Integration Formulas and the Net Change Theorem»; \cite{openstax2}, §2.1 «Areas between Curves».
\item Strang~\cite{strang}, §5.6 «Properties of the Integral and Average Value» y §5.7 «The Fundamental Theorem and Its Applications».
\item MIT 18.01SC~\cite{mit1801}, unidad~3, sesiones~47--52 (teorema fundamental y su demostración) y~60--61 (valor medio).
\end{lecturas}

\section*{Ejercicios del tema 2}
\addcontentsline{toc}{section}{Ejercicios del tema 2}

\begin{ejercicio}{1}{eje:2-conceptos}
Sea $F(x)=\int_0^{x}\sqrt{1+t^{3}}\dif t$ para $x\ge0$. (a) ¿Cuánto vale $F(0)$? (b) Calcule $F'(x)$ y $F''(x)$ indicando qué resultado usa. (c) ¿Es $F$ creciente? ¿Es convexa?
\end{ejercicio}

\begin{ejercicio}{1}{eje:2-error-barrow}
Explique el error en el cálculo $\int_{-1}^{1}\frac{\dif x}{x^{2}}=\eval{-\frac1x}{-1}{1}=-1-1=-2$. ¿Qué hipótesis de la regla de Barrow falla? ¿Por qué el signo del resultado ya delata el error?
\end{ejercicio}

\begin{ejercicio}{2}{eje:2-barrow}
Calcule con la regla de Barrow: (a) $\int_1^4\bigl(3\sqrt x-\frac{2}{x^{2}}\bigr)\dif x$; (b) $\int_0^{\pi/2}\cos x\dif x$; (c) $\int_0^1e^{2x}\dif x$; (d) $\int_0^1\frac{\dif x}{1+x}$; (e) $\int_{-1}^{2}\abs{x}\dif x$.
\end{ejercicio}

\begin{ejercicio}{2}{eje:2-arco-seno}
Halle el área bajo un arco de $y=\sin x$ y bajo un arco de $y=\sin(ax)$, con $a>0$ constante. ¿Cómo cambia el área al comprimir horizontalmente la gráfica?
\fuente{adaptado de MIT 18.01SC~\cite{mit1801}, \emph{18.01 Exercises}, ejercicio 3C-5.}
\end{ejercicio}

\begin{ejercicio}{2}{eje:2-parabola-eje}
(a) Halle el área de la región limitada por el eje $x$ y la curva $y=x^{2}-4$, y después por el eje $x$ y $y=x^{2}-a$, con $a>0$. (b) Halle el área de la región limitada por $y=x^{2}$ y $y=2x+3$, y la de la limitada por $y=\sqrt x$ y $y=x^{2}$.
\fuente{adaptado de MIT 18.01SC~\cite{mit1801}, \emph{18.01 Exercises}, ejercicio 3C-6.}
\end{ejercicio}

\begin{ejercicio}{2}{eje:2-derivar}
Calcule $F'(x)$ si: (a) $F(x)=\int_0^{x^{2}}\sqrt{u}\,\sin u\dif u$; (b) $F(x)=\int_0^{\sin x}\frac{\dif t}{\sqrt{1-t^{2}}}$, para $x\in(-\pi/2,\pi/2)$; (c) $F(x)=\int_x^{x^{2}}e^{-t^{2}}\dif t$. En (b), identifique $F$.
\fuente{adaptado de MIT 18.01SC~\cite{mit1801}, \emph{18.01 Exercises}, ejercicio 3D-7.}
\end{ejercicio}

\begin{ejercicio}{2}{eje:2-hallar-f}
Sea $f$ continua. Halle $f(\pi/2)$ si: (a) $\int_0^{x}f(t)\dif t=2x(\sin x+1)$ para todo $x$; (b) $\int_0^{x/2}f(t)\dif t=2x(\sin x+1)$ para todo $x$.
\fuente{adaptado de MIT 18.01SC~\cite{mit1801}, \emph{18.01 Exercises}, ejercicio 3D-8.}
\end{ejercicio}

\begin{ejercicio}{2}{eje:2-valor-medio}
Calcule el valor medio de: (a) $f(x)=x^{2}$ en $[0,3]$; (b) $f(x)=e^{x}$ en $[0,\ln2]$. En cada caso, halle el punto $c$ que proporciona el \cref{teo:tvm-integral} y compruebe que está en el intervalo.
\end{ejercicio}

\begin{ejercicio}{2}{eje:2-estudio-F}
Sea $F(x)=\int_1^{x}\frac{\dif t}{1+t^{2}}$. Estudie: (a) su dominio y el signo de $F$; (b) los intervalos de crecimiento, los extremos relativos y los puntos de inflexión; (c) que $F(x)<1$ para todo $x>1$, sin usar la función arcotangente (compare el integrando con $1/t^{2}$); (d) su expresión con $\arctan$ y sus límites en $\pm\infty$.
\fuente{adaptado de MIT 18.01SC~\cite{mit1801}, \emph{18.01 Exercises}, ejercicio 3D-3.}
\end{ejercicio}

\begin{ejercicio}{2}{eje:2-limite}
Calcule $\displaystyle\lim_{h\to0}\frac1h\int_1^{1+h}\frac{\dif t}{\sqrt{1+t^{4}}}$ de dos formas: (a) con el \cref{teo:tvm-integral}; (b) como la derivada en $x=1$ de $F(x)=\int_0^{x}\frac{\dif t}{\sqrt{1+t^{4}}}$.
\fuente{adaptado de MIT 18.01SC~\cite{mit1801}, \emph{18.01 Exercises}, ejercicio 3D-5.}
\end{ejercicio}

\begin{ejercicio}{3}{eje:2-par-impar}
Sea $f$ continua en $\R$ y $F(x)=\int_0^{x}f(t)\dif t$. Demuestre que: (a) si $f$ es par, $F$ es impar; (b) si $f$ es impar, $F$ es par. (Indicación: derive $F(x)+F(-x)$ o $F(x)-F(-x)$.)
\fuente{adaptado de MIT 18.01SC~\cite{mit1801}, \emph{18.01 Exercises}, ejercicio 3E-5.}
\end{ejercicio}

\begin{ejercicio}{3}{eje:2-logaritmo}
Sea $L(x)=\int_1^{x}\frac{\dif t}{t}$ para $x>0$. Usando solo el \cref{teo:tfc1} y la \cref{prop:primitivas-constante}, demuestre que: (a) $L(1/a)=-L(a)$ para todo $a>0$; (b) $L(a^{n})=n\,L(a)$ para todo $n\in\N$, por inducción; (c) $L(x)\to+\infty$ cuando $x\to+\infty$. (Indicación para (c): $L(2)>0$.)
\fuente{adaptado de MIT 18.01SC~\cite{mit1801}, \emph{18.01 Exercises}, ejercicio 3E-1.}
\end{ejercicio}

\begin{ejercicio}{3}{eje:2-periodica}
Sea $f$ continua en $\R$ y periódica de periodo $T>0$. Demuestre que $\int_a^{a+T}f(x)\dif x$ no depende de $a$. Deduzca que $\int_a^{a+2\pi}\sin^{2}x\dif x=\pi$ para todo $a$, sabiendo que $\int_0^{2\pi}\sin^{2}x\dif x=\pi$.
\end{ejercicio}

\begin{ejercicio}{3}{eje:2-simetrica-f}
Sea $f$ continua en $[a,b]$ tal que $\int_a^{x}f=\int_x^{b}f$ para todo $x\in[a,b]$. Demuestre que $f(x)=0$ para todo $x\in[a,b]$.
\end{ejercicio}

\begin{ejercicio}{3}{eje:2-tvm-ponderado}
Demuestre el teorema del valor medio ponderado: si $f$ y $w$ son continuas en $[a,b]$ y $w\ge0$, existe $c\in[a,b]$ con $\int_a^bf(x)\,w(x)\dif x=f(c)\int_a^bw(x)\dif x$. Muestre con un ejemplo que la conclusión puede fallar si $w$ cambia de signo.
\end{ejercicio}
